Their initial branches are $A_+(s)=1$ for $1\le s\le3$ and
$A_-(s)=0$ for $s\le2$, and subsequent branches satisfy
\[
 A_+'(s)=\frac{A_-(s-1)}{s-1},\qquad
 A_-'(s)=\frac{A_+(s-1)}{s-1}.
\]
The equation for a branch involves only the preceding unit interval.
One therefore obtains rational interval enclosures successively on a
rational mesh, splitting at every integer breakpoint.  Lower and upper
integral bounds follow from the corresponding interval bounds for the
integrand.  Chord bounds and curvature bounds sharpen these enclosures
without changing their direction.  An arbitrary query is enclosed on
its entire support, including both adjoining cells when an endpoint
is a mesh knot.  The analogous construction for $s\omega(s)$ uses
$(s\omega(s))'=\omega(s-1)$ and its initial branch
$\omega(s)=1/s$ on $[1,2]$.

The chord integration can be implemented without numerical
quadrature.  On $[j/n,(j+1)/n]$ the integral of the affine
endpoint interpolant of a numerator $V$, divided by its argument,
is
\[
 \bigl((j+1)\log(1+1/j)-1\bigr)V(j/n)
 +\bigl(1-j\log(1+1/j)\bigr)V((j+1)/n).
\]
Both weights are positive because they are integrals of
nonnegative affine basis functions against $du/u$.
The monotonicities of $F$ and $f$ imply that $A_+$ is convex
and $A_-$ is concave on their nonconstant branches.
Their numerator curvature on a delayed cell is bounded
in absolute value by $1/(j/n-1)$.  Thus the integrated
chord error is at most
\[
 \frac{1}{12nj(j-n)}\qquad(j>n).
\]
This quantity is subtracted from a chord lower estimate
or added to a chord upper estimate as appropriate.
The weights themselves have alternating rational series
\[
 (j+1)\log(1+1/j)-1
  =\sum_{k\ge1}\frac{(-1)^{k+1}}{k(k+1)j^k},\qquad
 1-j\log(1+1/j)
  =\sum_{k\ge1}\frac{(-1)^{k+1}}{(k+1)j^k}.
\]
Their fourth and fifth partial sums enclose the exact
weights.  This supplies rational enclosures even where
direct subtraction of nearly equal logarithms would be
numerically inconvenient.
